\section[Cote 98 : une théorie axiomatique des normes, le modèle des G-ensembles]{Cote 98 : une théorie axiomatique des normes, le modèle des $G$-ensembles\protect\verifie{Folder98}}\label{sec:98}

La cote 98, \emph{Schémas de Hilbert. Normes : notes manuscrites (s.d.)}, que l'inventaire date \q{[à partir de 1966-1967]}, pose aux pages~10 et~11 une axiomatique des normes dans une catégorie à sommes finies : une classe $\mathcal N$ de morphismes stable par identités, composition, changement de base, sommes et quotients $(X^n_Y)/H$, $H \subset \mathfrak S_n$ ((i) à (v)) ; un degré multiplicatif, invariant par changement de base, additif, avec $\deg\bigl((X^n_Y)/H \to Y\bigr) = \mathrm{card}(E^n/H)$ si $\deg u = r$ et $E = \{1, \ldots, r\}$ ((i bis) à (v bis)) ; des sections canoniques $\mathrm{Can}_{X/Y} : Y \to \mathrm{Sym}^r_Y(X)$ avec $\mathrm{Can}_{X/X} = \mathrm{id}$, le diagramme de transitivité de la page~11, le changement de base et l'additivité ((i ter) à (iv ter)) (\lot{98}{1}{10}, \lot{98}{1}{11}). Dans l'exemple~I, les $G$-ensembles, $u$ est de degré $r$ si toutes ses fibres ont $r$ éléments, et $\mathrm{Can}_{X/Y}(y)$ est la classe des points de la fibre \q{rangés dans un ordre quelconque} ; la page dit les axiomes \q{immédiats}, et aucun exemple n'est vérifié pour (ii ter).

\begin{theoreme}[D'après les pages 10 à 12, exemple I\piste{98-axiomatic-norms-can-section}]
Les applications dont toutes les fibres sont finies de même cardinal $r$, munies de ce cardinal comme degré, satisfont (i) à (v) et (i bis) à (v bis) ; la fibre de $(X^n_Y)/H \to Y$ en $y$ est en bijection avec $E^n/H$. Avec $\mathrm{Sym}^r_Y(X)$ formé des couples $(y, m)$, $m$ multi-ensemble de $r$ points de $u^{-1}(y)$, et $\mathrm{Can}_{X/Y}(y)$ le multi-ensemble des points de la fibre, (i ter) à (iv ter) sont satisfaits, et $\mathrm{Can}$ est compatible à l'action de $G$ quand $u$ est équivariant.
\end{theoreme}

\noindent L'énoncé tient : voir \texttt{hasDeg\_comp}, \texttt{hasDeg\_quotPow}, \texttt{can\_comp}, \texttt{can\_pullback}, \texttt{can\_sum} et \texttt{can\_smul} dans \href{https://github.com/Commutator-IO/grothendieck-project/blob/main/lean/Grothendieck/Folder98.lean}{\texttt{Folder98.lean}}. Pour (ii ter), la fibre de $vu$ au-dessus de $z$ est la réunion disjointe des fibres de $u$ au-dessus des $s$ points de la fibre de $v$. Le groupe ne joue aucun rôle au-delà de l'équivariance de $\mathrm{Can}$. Ne sont pas démontrés : l'exemple~II (morphismes finis localement libres et norme de la page~2), l'exemple~III (revêtements et $\Phi(E) = X \times^\Gamma E$) et la transitivité des normes des pages~6 à~8.
